\chapter{Charcot--Marie--Tooth Disease}

\section*{Definition and Overview}
Charcot--Marie--Tooth disease (CMT) is the most common inherited disorder of the peripheral nerves. It is not a single disease but a group of genetically distinct conditions that damage the nerves carrying messages between the brain and spinal cord and the rest of the body. The result is progressive muscle weakness and wasting, together with loss of sensation, that most often affects the feet, lower legs, hands, and forearms. The condition is named after the three physicians who first described it in the late nineteenth century. Symptoms typically begin in childhood or early adulthood. Although there is no cure, most affected people remain mobile throughout life and have a normal life span.

\section*{Epidemiology}
Charcot--Marie--Tooth disease affects approximately 1 in 2,500 people worldwide, making it one of the most common inherited neurological conditions. It occurs in all ethnic groups, and males and females are affected in roughly equal numbers, although certain genetic subtypes show sex-specific patterns of inheritance. Because CMT is inherited, it often runs in families, yet in a substantial proportion of cases there is no family history and the causative mutation is new. Symptoms usually appear before the age of 20, but some severe forms present in infancy, while milder subtypes may go unrecognised until middle or later adulthood.

\section*{Aetiology and Pathophysiology}
\medskip
\textbf{Genetic basis.} CMT is caused by mutations in more than 100 identified genes that encode proteins involved in the structure or function of peripheral nerves:
\begin{itemize}
    \item \textbf{Inheritance patterns:} the disorder may be autosomal dominant, autosomal recessive, or X-linked, depending on the gene involved
    \item \textbf{Demyelinating forms:} the protective myelin sheath around nerves fails to form correctly or breaks down, slowing conduction of nerve signals
    \item \textbf{Axonal forms:} the nerve fibre itself degenerates, reducing the strength and integrity of nerve signals
    \item \textbf{Length-dependent damage:} the longest nerve fibres are affected first, which explains why symptoms begin in the feet and hands and progress towards the trunk
\end{itemize}

\medskip
\textbf{Pathophysiology.} Whatever the genetic cause, the final common pathway is dysfunction and degeneration of peripheral nerves. Muscles that are no longer adequately stimulated by healthy nerves gradually shrink and weaken. Damage to the sensory fibres reduces feeling, and damage to the autonomic fibres can produce coldness and circulatory changes in the extremities.

\section*{Clinical Features}
\begin{itemize}
    \item \textbf{Foot deformities:} high foot arches (pes cavus), and less often flat feet, together with hammertoes
    \item \textbf{Foot drop:} difficulty lifting the front of the foot, producing a characteristic high-stepping or slapping gait
    \item \textbf{Gait problems:} difficulty walking, frequent tripping, and repeated ankle sprains
    \item \textbf{Leg wasting:} weakness and wasting of the lower leg muscles, giving an "inverted champagne bottle" or stork-leg appearance
    \item \textbf{Hand symptoms:} weak grip, clumsiness, and difficulty with fine tasks such as writing and doing up buttons
    \item \textbf{Sensory symptoms:} reduced sensation, numbness, tingling, or pain in the hands and feet
    \item \textbf{Reflex changes:} reduced or absent tendon reflexes such as the knee jerk
    \item \textbf{Other features:} cold extremities, tremor, and in some subtypes mild hearing loss, breathing difficulty, or scoliosis
    \item \textbf{Course:} symptoms are usually symmetrical, slowly progressive, and highly variable in severity
\end{itemize}

\section*{Differential Diagnosis}
\begin{itemize}
    \item Chronic inflammatory demyelinating polyneuropathy (CIDP)
    \item Hereditary neuropathy with liability to pressure palsies (HNPP)
    \item Metabolic and nutritional neuropathies, including diabetic and vitamin B12 deficiency neuropathy
    \item Spinal muscular atrophy and other motor neuron diseases
    \item Friedreich's ataxia and other inherited ataxias
    \item Guillain--Barr\'e syndrome in acute presentations
    \item Muscular dystrophy and other primary muscle diseases
    \item Alcohol-related peripheral neuropathy
    \item Poliomyelitis or post-polio syndrome
\end{itemize}

\section*{When to Seek Medical Care}
See a doctor if you or your child notice weakness, wasting, numbness, or tingling in the feet or hands, foot drop, a high-stepping walk, frequent tripping, or progressive difficulty with walking and hand function. Anyone with a family history of CMT who develops such symptoms should also be assessed.

\medskip
\textbf{Call 000 (or local emergency services) for:}
\begin{itemize}
    \item Sudden weakness or paralysis, especially of recent onset
    \item Sudden difficulty breathing or swallowing
    \item Rapid worsening of symptoms over days or weeks
    \item Loss of consciousness or a seizure
\end{itemize}

\section*{Diagnosis and Investigations}
\begin{itemize}
    \item \textbf{History:} family history of similar symptoms, age of onset, and the pattern and progression of weakness and sensory loss
    \item \textbf{Neurological examination:} assessment of muscle strength and bulk, reflexes, coordination, gait, and sensation in all limbs
    \item \textbf{Nerve conduction studies:} measure how quickly electrical signals travel along peripheral nerves and distinguish demyelinating from axonal forms
    \item \textbf{Electromyography (EMG):} examines the electrical activity of muscle to confirm denervation
    \item \textbf{Genetic testing:} a blood test for mutations in known CMT genes, which confirms the diagnosis and informs genetic counselling in many cases
    \item \textbf{Nerve biopsy:} now rarely required, since genetic testing has become widely available
    \item \textbf{Specialist referral:} to a neurologist and, where available, a neuromuscular or CMT genetics clinic
\end{itemize}

\section*{Management}
\begin{itemize}
    \item \textbf{Physiotherapy:} exercises to maintain strength, flexibility, and balance, and to prevent contractures and falls
    \item \textbf{Occupational therapy:} help with hand function, everyday activities, and the provision of adaptive equipment
    \item \textbf{Orthotics:} ankle-foot orthoses and insoles to correct foot drop and improve gait
    \item \textbf{Footwear:} well-fitted, supportive shoes, sometimes custom-made
    \item \textbf{Surgery:} foot and ankle procedures to correct painful deformities such as high arches and hammertoes, or to release contractures
    \item \textbf{Pain management:} medicines and other strategies for neuropathic pain where it occurs
    \item \textbf{Hand therapy:} splints, exercises, and ergonomic aids to preserve grip and fine motor skills
    \item \textbf{Genetic counselling:} to discuss inheritance, implications for relatives, and family planning
    \item \textbf{Support services:} disability supports, mobility aids, and peer support organisations
    \item \textbf{Follow-up:} regular review by a neurologist and rehabilitation team so management can be adjusted as the condition evolves
\end{itemize}

\section*{Complications}
\begin{itemize}
    \item Progressive weakness and muscle wasting in the limbs
    \item Foot drop and an unsteady, high-stepping gait
    \item Falls, ankle sprains, and fractures
    \item Joint contractures from fixed tightening of tendons and ligaments
    \item Curvature of the spine (scoliosis)
    \item Chronic neuropathic pain, burning, or numbness in the extremities
    \item Reduced hand function affecting writing, dressing, and employment
    \item Pressure injuries on the feet from reduced sensation
    \item In some subtypes, breathing difficulties, hearing loss, or optic nerve involvement
    \item Fatigue and reduced quality of life
    \item Psychological impact, including anxiety and low mood
\end{itemize}

\section*{Prognosis and Outcomes}
Charcot--Marie--Tooth disease is slowly progressive, and in most people the rate of deterioration is gradual. It does not generally shorten life expectancy, and most affected individuals remain able to walk, although walking may become more effortful with age and many people eventually use ankle supports or walking aids. Severity varies widely, even within the same family: some people have very mild symptoms that never cause significant disability, while others require a wheelchair. There is currently no treatment that halts or reverses the nerve damage, but physiotherapy, orthotics, and surgery substantially improve function and quality of life. The outlook depends on the genetic subtype and should be discussed with a neurologist.

\section*{Prevention and Self-Care}
\begin{itemize}
    \item Stay physically active with low-impact exercise such as swimming, cycling, and stretching
    \item Follow a structured physiotherapy program to maintain strength and flexibility
    \item Wear prescribed splints, braces, and supportive footwear as recommended
    \item Inspect the feet daily for cuts, blisters, and pressure sores, since sensation may be reduced
    \item Reduce fall risk by keeping the home well lit and free of trip hazards
    \item Maintain a healthy weight to lessen strain on weakened muscles and joints
    \item Avoid heavy alcohol use, which can further damage peripheral nerves
    \item Seek genetic counselling before starting a family if CMT runs in the family
    \item Use ergonomic aids for writing, typing, and gripping
    \item Connect with CMT support organisations for information, advocacy, and peer support
\end{itemize}

\section*{Sources and Further Reading}
The following reputable sources provide further information on Charcot--Marie--Tooth disease:
\begin{itemize}
    \item NHS. \textit{Health A--Z: conditions affecting the nervous system.} https://www.nhs.uk/conditions/
    \item Mayo Clinic. \textit{Diseases and conditions library.} https://www.mayoclinic.org/
    \item MSD Manual. \textit{Professional edition: peripheral nerve disorders.} https://www.msdmanuals.com/
    \item MedlinePlus. \textit{Health topics on genetic and neurological conditions.} https://medlineplus.gov/
    \item Centers for Disease Control and Prevention. \textit{Genetic and neurological conditions resources.} https://www.cdc.gov/
\end{itemize}
